\section{Streams and Clocks}
\label{sec:semantics:streams}

In dataflow languages like \GRust, an expression is interpreted not as a single value, but as an
\emph{infinite sequence of values} it produces over time. This sequence is called a \emph{stream}, a
concept foundational to the semantics of dataflow programming~\cite{%
    DBLP:conf/ifip/Kahn74,%
    DBLP:journals/tecs/BourkePP23,%
    DBLP:journals/tcad/LeeS98%
}. Streams provide a mathematical model for the entire temporal evolution of a value.
%
This section formalizes the theory of streams and describes its application to the denotational
semantics of \GRust.

\subsection{The Principle of Coinduction}
\label{sec:semantics:streams:coinduction}

Reasoning about infinite data structures requires a different approach from the induction used for
finite data structures, such as lists. Inductive proofs rely on a base case (\eg an empty list) and
a step that decomposes a structure into a \emph{smaller} substructure (\eg the head and tail of a
list). This process is guaranteed to terminate because the structure diminishes at each step.

Infinite streams, by definition, lack a base case and cannot be reduced to smaller parts. Attempting
an inductive proof on a stream leads to infinite regress: proving a property for a stream would
require proving it for its tail, which is also an infinite stream. The proof would never conclude.

Instead, we use the dual principle of \emph{coinduction}, a technique for defining and proving
properties of infinite data structures~\cite{DBLP:journals/mscs/KozenS17}. Coinduction is based on
the idea of \emph{observation}: what can be concretely observed about an infinite object in a finite
amount of time? For a stream, we can only observe its first element (the \emph{head}) and the
remainder of the stream (the \emph{tail}).

Coinduction formalizes this observational logic into a proof method. To prove that a property holds
for an infinite stream, we must establish two points:
%
\begin{description}
    \item[Progress] The proof must show that the property holds for the head of the stream---the
          part that can be observed immediately. This step ensures that the stream is productive and
          the proof is not vacuous.
    \item[Opacity] After demonstrating progress, the proof may \emph{assume} that the property holds
          for the tail of the stream. This assumption is the \emph{coinductive hypothesis}. The term
          \emph{opacity} signifies that we treat the rest of the stream as an abstract entity that
          satisfies the hypothesis, without inspecting its infinite structure further.
\end{description}
%
By requiring observable progress before applying the hypothesis, coinduction avoids infinite regress
and provides a sound method for verifying properties of infinite data structures.

\subsection{Definition of Streams}
\label{sec:semantics:streams:def}

Let $A$ be a non-empty set. A stream over $A$, or an $A$-stream, is an infinite sequence of elements
from $A$. The set of all $A$-streams is denoted by \Stream{A}.

Formally, a stream is defined coinductively as the greatest fixed-point of the recursive type
equation:
%
\begin{equation}
    \Stream{A} = \nu X. A \times X \label{eq:semantics:streams:type_equation}
\end{equation}
%
This equation states that the type \Stream{A} is equivalent to a pair, consisting of an element from
$A$ and another stream. This structure implies functions to construct and deconstruct streams.

The unique constructor for streams is the function $\texttt{cons}$:
%
\begin{align}
    \texttt{cons}: A \times \Stream{A} \to \Stream{A}
\end{align}
%
For readability, we adopt the infix operator \step as a convenient notation for this constructor:
for any $a \in A$ and $s \in \Stream{A}$, the stream $a \step s$ denotes $\texttt{cons}(a, \, s)$.
Using this notation, we can define specific streams with recursive equations. For example, a stream
of zeros is the unique solution to:
%
\begin{equation}
    \texttt{zeros} = 0 \step \texttt{zeros}
\end{equation}

Conversely, two observer functions deconstruct a stream. The function \texttt{head} returns the
first element of the stream:
\begin{align}
     & \texttt{head} : \Stream{A} \to A \\
     & \texttt{head}(a \step s) = a
\end{align}
%
The function $\texttt{tail}$ returns the remainder of the stream:
\begin{align}
     & \texttt{tail} : \Stream{A} \to \Stream{A} \\
     & \texttt{tail}(a \step s) = s
\end{align}
%
The interaction between the constructor and observers is governed by the \emph{unfold property} that
follows. For any stream $s \in \Stream{A}$:
%
\begin{align}
    s = \texttt{head}(s) \step \texttt{tail}(s) \label{eq:semantics:streams:unfold}
\end{align}
%
The unfold property asserts that any stream can be deconstructed and then perfectly reconstructed.

The equivalence of two $A$-streams, $s_1$ and $s_2$, asserts they are pointwise identical and is
written $s_1 \equiv s_2$. We define this relation using the coinductive principle, formalized by the
inference rule \nameref{eq:semantics:streams:equiv}.
%
\begin{mathpar}
    \myinfer
    {a_1 = a_2 \and s_1 \equiv s_2}
    {(a_1 \step s_1) \equiv (a_2 \step s_2)}
    {Equiv}
    \label{eq:semantics:streams:equiv}
\end{mathpar}
%
This rule is a direct application of coinductive reasoning. To prove two streams equal, we must
satisfy the two conditions in the premise:
%
\begin{enumerate}
    \item \textbf{Progress:} Show that their heads are equal ($a_1 = a_2$).
    \item \textbf{Opacity:} Assume, \via the coinductive hypothesis, that their tails are also
          equivalent ($s_1 \equiv s_2$).
\end{enumerate}

The stream formalism is the foundation of our denotational semantics. The constructor \step is used
to define stream transformations, while the equivalence relation $\equiv$ is central to semantic
properties. Although the \texttt{head} and \texttt{tail} functions are theoretically fundamental,
they do not appear explicitly in our semantic rules. Instead, we use pattern matching on the form
$a \step s$ to deconstruct a stream into its head $a$ and tail $s$.

\subsection{Application in Denotational Semantics}
\label{sec:semantics:streams:usage}

In the denotational semantics of \GRust, streams provide the mathematical basis for modeling
dataflows. The meaning of a program is given by associating each identifier with an infinite stream
of values. The structure of these streams differs between \GRSync and \GRFRP to reflect their
distinct execution models.

\paragraph{\GRSync: Synchronous Streams}

The semantics of \GRSync follows the synchronous dataflow paradigm~\cite{%
    DBLP:journals/tecs/BourkePP23,%
    DBLP:conf/ifip/Kahn74%
}. In this model, program execution advances in lock-step through a sequence of discrete
\emph{logical instants}. Each identifier is interpreted as a stream, where the $n^\text{th}$ element
corresponds to the value of that identifier at the $n^\text{th}$ logical instant. This synchronous
approach ensures that values at the same position across different streams are temporally aligned
and thus available simultaneously during execution.

A challenge, however, arises with conditional computations, such as expressions in a specific branch
of a pattern match. These expressions produce values only when their branch is active. To preserve
the lock-step nature of streams, every stream must contain an entry for every logical instant.
%
This issue is resolved by introducing a special marker for absence. Let \Val be the set of all
possible values. We extend this set to $\ValBot = \Val \cup \{\bot\}$, where the distinct
\emph{bottom} value, $\bot$, represents the absence of a value at a given instant. This ensures that
even when a computation is inactive, its corresponding stream is populated with a placeholder,
maintaining alignment with all other streams.

The denotational semantics of \GRSync therefore maps each identifier to a stream of type
\Stream{\ValBot}, and components to functions from input streams to output streams.

\paragraph{\GRFRP: Asynchronous Timestamped Streams}

In contrast to \GRSync, \GRFRP models asynchronous services by abandoning the notion of a global
logical instant. This design reflects the reality of reactive systems where inputs from sensors or
other components arrive at irregular intervals. To manage this asynchrony, we must establish a
coherent temporal ordering of values.
%
The solution is to pair each value with a timestamp, a model similar to the
\emph{Tagged Signal Model}~\cite{DBLP:journals/tcad/LeeS98}. Let \Time be the set of all possible
timestamps. A flow in \GRFRP is interpreted as a stream where each element is a pair $(t, v)$,
consisting of a timestamp $t \in \Time$ and a value $v \in \Val$.

A critical property of these streams is that time must not flow backward. We enforce this by
requiring timestamps to be monotonically non-decreasing. A stream that respects this ordering is a
well-formed \emph{timestamped stream}. The set of all such streams over a value set $A$, denoted by
\TStream{A}, is the subset of \Stream{\Time \times A} where timestamps are ordered.
%
Formally, this property is defined coinductively:
%
\begin{equation}
    \TStream{A} = \{ (t, v) \step (t', v') \step \ts \mid t \le t' \wedge
    ((t', v') \step \ts) \in \TStream{A}\}
    \label{eq:semantics:streams:usage:tstream}
\end{equation}

The denotational semantics of \GRFRP interprets flows as well-formed timestamped streams of type
\TStream{\Val}, and services as functions from input to output timestamped stream contexts. This
model captures the asynchronous arrival of data while preserving a well-defined progression of time.

\subsection{Clocks}
\label{sec:semantics:streams:clock}

In synchronous and reactive languages, a \emph{clock} is a stream that governs the rate of other
streams~\cite{DBLP:conf/emsoft/ColacoP03}. It dictates when computations are performed and when
values are considered present. This concept is essential for coordinating dataflows, particularly
those operating at different rates. Clocks are handled differently in \GRSync and \GRFRP to suit
their respective execution models.

\paragraph{Implicit Clocks in \GRSync}

In the synchronous context of \GRSync, a clock, denoted by $ck$, is a stream of booleans,
\Stream{\Bool}. Each element corresponds to a logical instant. If the value is \true, the clock is
said to \emph{tick}, signaling that associated computations should execute. If the value is \false,
computations on that clock are inactive.

Every stream is defined \emph{on} a clock, which determines the presence of its values. We write
\on{s}{ck} to assert that stream $s$ is on clock $ck$. This relation means that at any instant, if
the clock ticks \true, the stream must contain a value from \Val; if the clock is \false, the stream
must contain the bottom value, $\bot$. This is defined by the following two rules.

\begin{mathpar}
    \myinfer
    {\on{s}{ck}}
    {\on{v \step s}{\true \step ck}}
    {On{\scriptsize val}}
    \label{rule:semantics:streams:onclock:val}

    \myinfer
    {\on{s}{ck}}
    {\on{\bot \step s}{\false \step ck}}
    {On{\scriptsize bot}}
    \label{rule:semantics:streams:onclock:bot}
\end{mathpar}

A key design choice in \GRSync is that clocks are \emph{implicit}. Programmers do not manipulate
them directly; instead, the compiler generates sub-clocks automatically from the program structure.
For instance, each branch of a \kwf{match} expression defines a sub-clock that ticks only when
the input value matches the corresponding pattern. Expressions within that branch are computed on
this sub-clock. If the patterns are exhaustive, the outputs from all branches are merged back into a
single stream on the original clock of the matched expression.

This approach contrasts with languages that feature \emph{explicit} clock management, which require
a \emph{clock calculus}~\cite{DBLP:conf/emsoft/ColacoP03} to verify that clock relations are
correctly defined by the programmer. While powerful, mastering a clock calculus demands a deep
understanding of the formal semantics of the language. Indeed, many of these languages provide
higher-level constructs, such as \kwf{automaton} or \kwf{switch}, to ease the development of
state-based systems. Following this principle, \GRSync omits entirely explicit clocks and instead
provides the higher-level constructs \kwf{when} and \kwf{match}. This design ensures that programs
are well-clocked by construction, a property formally established by the productivity theorem
(\Cref{thm:semantics:grsync:theorems:productivity}).

All top-level computations within a \GRSync component are defined on a \emph{base clock}. Formally,
this clock is an infinite stream of \true values: $ck = \true \step ck$.
%
This clock provides the fundamental sequence of logical instants that underpins the synchronous
execution model, representing the fastest possible rate within the top-level component. Every
stream in the top-level component is on this clock, meaning it must provide a value at each instant.
%
The base clock is an abstraction that connects the synchronous semantics of \GRSync to the
asynchronous environment of \GRust. This connection is realized by the calling \GRFRP service, which
acts as the driver for the \GRSync component, following the
\emph{Globally Asynchronous, Locally Synchronous} (GALS) principle~\cite{DBLP:phd/us/Chapiro85}.
%
The \GRFRP service orchestrates the execution. When an asynchronous event occurs (\eg a new sensor
value arrives), the service invokes the top-level component to process the new data. Each such
invocation corresponds to a single logical instant, or one \true tick, of the base clock.
%
Therefore, the sequence of invocations from the \GRFRP service constitutes the concrete realization
of the abstract base clock. From the perspective of the component, its execution is a deterministic
progression from one discrete instant to the next, unaware of the real-time delays or the
asynchronous nature of its driver. This design effectively bridges the two domains: the \GRFRP
service translates unpredictable external events into a predictable sequence of synchronous ticks,
thereby ensuring deterministic computation within a globally asynchronous system.

\paragraph{Timestamped Clocks in \GRFRP}

In the asynchronous setting of \GRFRP, the lock-step model is replaced by explicit timestamps.
Consequently, a clock is a well-formed \emph{timestamped stream} of booleans, denoted by
$\tck \in \TStream{\Bool}$ (following the definition of \TStream{A} from
\Cref{eq:semantics:streams:usage:tstream}, with $A$ being \Bool). A tick of this clock is a pair
$(t, \true)$, indicating that a computation is active at time $t$.

This model defines a formal relationship between a dataflow and its controlling clock. We say a
stream \ts is \emph{on} a clock \tck, written \on{\ts}{\tck}, if the timestamps of the values in \ts
are a subset of the timestamps where \tck ticks \true. This property ensures a stream carries values
only at times permitted by its clock. The relation is defined coinductively by the following two
rules:

\begin{mathpar}
    \myinfer
    {\on{\ts}{\tck}}
    {\on{(t, v) \step \ts}{(t, \true) \step \tck}}
    {On{\scriptsize exact}}
    \label{rule:semantics:streams:onclock:exact}

    \myinfer
    {\on{\ts}{\tck}}
    {\on{\ts}{(t, b) \step \tck}}
    {On{\scriptsize pass}}
    \label{rule:semantics:streams:onclock:pass}
\end{mathpar}

The \nameref{rule:semantics:streams:onclock:exact} rule enforces direct synchronization: if a stream
has a value $v$ at time $t$, its clock must have a corresponding tick $(t, \true)$. The remainder of
the stream must then be on the remainder of the clock.

In contrast, the \nameref{rule:semantics:streams:onclock:pass} rule accommodates asynchrony. It
allows the clock to have a tick $(t, b)$ at a time when the data stream has no value. In this case,
the relation continues to be checked on the remainder of the clock. This rule permits a stream to be
sparser than its clock, which is fundamental to modeling asynchronous and event-driven behaviors.

\medskip

In summary, the concept of a clock in \GRust is adapted to the specific needs of its two
sub-languages. In the synchronous kernel, \GRSync, clocks are an implicit and structural mechanism,
ensuring lock-step execution and temporal alignment of data. This design guarantees that programs
are well-clocked by construction.
%
Conversely, in the asynchronous kernel, \GRFRP, timestamped clocks provide an explicit
representation of time, enabling the coordination of event-driven and temporally irregular
dataflows.
%
This dual treatment is central to the design of \GRust, as it allows the seamless integration of
deterministic, synchronous components within a flexible, asynchronous architecture. With these
foundational concepts of streams and clocks established, we are now equipped to define formally the
\GRust language.

\subsection{Advanced Coinductive Constructs}
\label{sec:semantics:streams:coinductive_constructs}

The denotational semantics of \GReact operations, detailed in \Cref{sec:semantics:grfrp:denot}, are
built upon auxiliary functions that generate timestamped streams coinductively. Proving that these
functions are well-defined is crucial for the soundness of the semantics. This section lays the
groundwork for these proofs by introducing and validating key coinductive constructs, such as
periodic time streams and stream pairs. These formal tools are essential for establishing the
denotational semantics of \GRFRP.

\paragraph{Periodic and Translated Time Streams}

A common requirement is to model periodic events, such as those from a timer. We can define a
periodic time stream coinductively. Let $T \in \Time$ be a period. The stream of periodic timestamps
$S_0^T$ is the unique solution to the equation:
%
\begin{equation}
    S_k^T = kT \step S_{k+1}^T
\end{equation}
%
This equation defines an infinite stream of timestamps $S_0^T = 0 \step T \step 2T \step \dots$.

More generally, we need to transform existing time streams. The function \texttt{trans} performs a
translation of a time stream to a new starting time while preserving its internal rhythm (\ie, the
intervals between its elements). It is defined as follows:
%
\begin{align}
     & \texttt{trans} : \Time \times \Stream{\Time} \to \Stream{\Time}                                       \\
     & \texttt{trans}(t, \: t_1 \step t_2 \step s) = t \step \texttt{trans}(t + (t_2 - t_1), \: t_2 \step s)
\end{align}
%
This function is well-defined by coinduction for its domain. It makes a recursive call on a stream
that has progressed, using the tail of the original input. By the coinductive hypothesis, we assume
this recursive call produces a valid stream. Thus, the entire operation yields a valid stream.
%
As a convenient shorthand, we write $t_0 + S_k^T$ to denote the translation of a periodic stream,
\ie $\texttt{trans}(t_0, S_k^T)$.

\paragraph{Stream Pairs}

To handle functions with multiple inputs or outputs, it is useful to define a pair of streams as a
single coinductive data structure. The set of all stream pairs over sets $A$ and $B$ is denoted by
\PStream{A}{B}. Formally, it is the greatest fixed-point of the recursive type equation:
%
\begin{equation}
    \PStream{A}{B} = \nu (X,Y). (A \times X) \times (B \times Y)
\end{equation}
%
We introduce constructors and observers to operate on each stream within the pair independently. The
constructors $\texttt{cons}_1$ and $\texttt{cons}_2$ prepend an element to the first or second
stream, respectively:
%
\begin{align}
    \texttt{cons}_1 & : A \times \PStream{A}{B} \to \PStream{A}{B} \\
    \texttt{cons}_1 & (a, (s_1, s_2)) = (a \step s_1, s_2)         \\[1em]
    \texttt{cons}_2 & : B \times \PStream{A}{B} \to \PStream{A}{B} \\
    \texttt{cons}_2 & (b, (s_1, s_2)) = (s_1, b \step s_2)
\end{align}
%
Conversely, the observers extract the head or tail from a specified stream. For the first stream in
the pair:
%
\begin{align}
    \texttt{head}_1 & : \PStream{A}{B} \to A   & \texttt{tail}_1 & : \PStream{A}{B} \to \PStream{A}{B} \\
    \texttt{head}_1 & ((a \step s_1, s_2)) = a & \texttt{tail}_1 & ((a \step s_1, s_2)) = (s_1, s_2)
\end{align}
%
Symmetrical functions, $\texttt{head}_2$ and $\texttt{tail}_2$, exist for the second stream:
%
\begin{align}
    \texttt{head}_2 & : \PStream{A}{B} \to B   & \texttt{tail}_2 & : \PStream{A}{B} \to \PStream{A}{B} \\
    \texttt{head}_2 & ((s_1, b \step s_2)) = b & \texttt{tail}_2 & ((s_1, b \step s_2)) = (s_1, s_2)
\end{align}
%
This abstraction allows functions to consume from or produce to one stream in a pair while leaving
the other unmodified.

For example, merging two ordered streams can be defined coinductively as follows:
%
\begin{align*}
    \texttt{merge}(a \step s_1, b \step s_2) & = a \step \texttt{merge}(s_1, b \step s_2) &
                                             & \text{if } a < b                             \\
    \texttt{merge}(a \step s_1, b \step s_2) & = b \step \texttt{merge}(a \step s_1, s_2) &
                                             & \text{otherwise}
\end{align*}
%
This function is well-defined by coinduction. The argument proceeds as follows. At each step, the
function deconstructs the input stream pair to compare the heads of its elements. Based on the
comparison, it makes a recursive call on a new pair formed from the tail of one of the original
streams (applying $\texttt{tail}_1$ or $\texttt{tail}_2$). This demonstrates \emph{progress}, as the
state of the coinductive pair has advanced.
%
Because progress is established, we can apply the \emph{coinductive hypothesis}. This allows us to
assume that the result of the recursive call is a valid, infinite stream. This step embodies the
principle of \emph{opacity}: we rely on the hypothesis for the remainder of the computation without
needing to inspect the infinite structure of the tails.
%
The function then uses the stream constructor to prepend the chosen head element ($a$ or $b$) to the
valid stream returned by the recursive call. Since this construction always yields a head followed
by a valid tail (by hypothesis), we conclude that the \texttt{merge} function is well-defined.

The \texttt{merge} function serves as a key example illustrating a crucial property of these
coinductive stream pairs: a proof can be progressed by advancing just one of the streams.

\subsection{Summary}
This section has laid the mathematical groundwork for our semantic definitions. We have formalized
the concept of infinite dataflows using the coinductive theory of streams and established two
distinct clocking models: implicit clocks for the synchronous \GRSync kernel and explicit
timestamped clocks for the asynchronous \GRFRP kernel. With these foundational concepts in place, we
are now equipped to formally define the \GRust language. We begin with \GRSync, the synchronous
component-modeling kernel, detailing its syntax and semantics.
